\chapter{Topology meets algebra}\label{chap:algebra}

In this chapter we show that the fundamental group construction respects the maps between topological spaces\footnote{In a more fancy language, it is \index{functor}\emph{functorial}.} and behaves predictably when one changes the base point of the space.
A bit later, we will calculate the fundamental group of the circle --- this will be the first example of a space with a nontrivial fundamental group.
Once this group is calculated, the results of this chapter will lead to applications.

\section{Induced homomorphism}\label{sec:Induced homomorphism}

A \index{pointed space}\emph{pointed space} (briefly $(\c{X},x_0)$) is a topological space $\c{X}$ with a choice of base point $x_0\in \c{X}$;
so fundamental group is a construction that produces a group $\pi_1(\c{X},x_0)$ from a pointed space $(\c{X},x_0)$.

A \index{pointed map}\emph{pointed map} (briefly $f\:(\c{X},x_0)\to (\c{Y},y_0)$) is a map $f\:\c{X}\z\to \c{Y}$ with a choice of base points $x_0\in \c{X}$ and $y_0\in \c{Y}$ such that $y_0=f(x_0)$.

Let $f\:(\c{X},x_0)\to (\c{Y},y_0)$ be a continuous pointed map.
Suppose that $a_0$ and $a_1$ are loops based at $x_0$ in $\c{X}$.
Then $(f\circ a_0)$ and $(f\circ a_1)$ are loops based at $y_0$ in $\c{Y}$.
Moreover,
\[f\circ(a_0*a_1)=(f\circ a_0)*(f\circ a_1)\]
Indeed, applying the definition of the product of paths and composition of maps to $f\circ(a_0*a_1)(t)$ and $(f\circ a_0)*(f\circ a_1)(t)$ we get exactly the same expression:
\begin{align*}
\begin{cases}
f\circ a_0(2\cdot t)&\text{if}\ t\le \tfrac12,
\\
f\circ a_1(2\cdot t-1)&\text{if}\ t\ge \tfrac12.
\end{cases}
\end{align*}

Further, note that
\[a_0\sim a_1 \quad\Rightarrow\quad f\circ a_0\sim f\circ a_1.\]
Indeed, if $a_\tau$ is a homotopy from $a_0$ to $a_1$, then $f\circ a_\tau$ is a homotopy from $f\circ a_0$ to $f\circ a_1$.

It follows that the map $f\:(\c{X},x_0)\to (\c{Y},y_0)$ \index{induced homomorphism}\emph{induces} a homomorphism
\[f_*\:\pi_1(\c{X},x_0)\to \pi_1(\c{Y},y_0)\]
that is defined by $f_*\:[a]\mapsto [f\circ a]$.
(The homomorphism $f_*$ maps the homotopy class of loop $a$ to the homotopy class of loop $f\circ a$.)

Evidently, the identity map of $(\c{X},x_0)$ induces the identity homomorphism on $\pi_1(\c{X},x_0)$.
Furthermore, for any continuous pointed maps
\[(\c{X},x_0)\stackrel{f}{\to}(\c{Y},y_0)\stackrel{g}{\to}(\c{Z},z_0),\]
we have $g_*\circ f_*=(g\circ f)_*$.
These two observations will be extremely useful.

\begin{thm}{Exercise}\label{ex:retract*}
Let $A$ be a retract of $\c{X}$ with retraction $r\:\c{X}\z\to A$.
Choose $x_0\in A$; let $i\:A\to\c{X}$ be the inclusion map.
Show the following.
\begin{subthm}{ex:retract*:0}
$r_*\:\pi_1(\c{X},x_0)\to \pi_1(A,x_0)$ is a surjective homomorphism.
\end{subthm}

\begin{subthm}{ex:retract*:i}
$i_*\:\pi_1(A,x_0)\to \pi_1(\c{X},x_0)$ is an injective homomorphism.
\end{subthm}

\begin{subthm}{ex:retract*:ri}
If $A$ is a weak deformation retract, then $r_*$ and $i_*$ are isomorphisms.
\end{subthm}


\end{thm}

\section[About notation f⁎]{About notation $\bm{f_*}$}

The induced  homomorphism \index{$f_*$}$f_*\:\pi_1(\c{X},x_0)\to \pi_1(\c{Y},y_0)$ has the same direction as the original map $f\:(\c{X},x_0)\to (\c{Y},y_0)$, by that reason it is also called \index{pushforward}\emph{pushforward}.
The symbol $f_*$ (read ``$f$-star'') can be used for \textit{any} induced map of this type; for example,
one can write $f_*$ for the induced map from the set of path-connected components of $\c{X}$ to the set of path-connected components of $\c{Y}$.
The associated constructions (fundamental group or set of path-connected components) are called \index{covariant}\emph{covariant}.

If you see $f_*$, then it means it is induced by $f$, goes in the same direction and you have to guess from the context the source and target of $f_*$.

For some other constructions, the induced map might go in the opposite direction.
In this case, the induced map of $f$ is denoted by $f^*$, and called \index{pullback}\emph{pullback}.
The associated constructions are called \index{contravariant}\emph{contravariant}.

As an example of contravariant construction, consider the so-called \emph{first cohomotopy group} $\pi^1(\c{X},x_0)$, which is the set of homotopy classes of pointed continuous maps $(\c{X},x_0)\to (\SS^1,s_0)$.
(The set $\pi^1(\c{X},x_0)$ has a natural structure of abelian group, but this is not important for us.)
Then any pointed continuous map $f\:(\c{X},x_0)\to (\c{Y},y_0)$ has a pullback $f^*\:\pi^1(\c{Y},y_0)\to\pi^1(\c{X},x_0)$.
Indeed, for a map $h\:(\c{Y},y_0)\to (\SS^1,s_0)$, the composition $h\circ f$ provides a good choice of map $(\c{X},x_0)\to (\SS^1,s_0)$, so we can define $f^*[h]\df [h\circ f]$.
On the other hand, there is no good choice of map $(\c{Y},y_0)\to (\SS^1,s_0)$ for a given map $(\c{X},x_0)\to (\SS^1,s_0)$.


\section{Transport of base point}

The following theorem provides a more exact version of \ref{prop:sc=1hclass}

{

\begin{wrapfigure}[5]{r}{27 mm}
\vskip-5mm
\centering
\includegraphics{mppics/pic-1415}
\end{wrapfigure}

\begin{thm}{Theorem}
Let $c$ be a path from $x_0$ to $x_1$ in a topological space $\c{X}$.
Then $[a]\mapsto [\bar c*a*c]$ defines an isomorphism of fundamental groups $\phi_c\:\pi_1(\c{X},x_0)\to \pi_1(\c{X},x_1)$.
\end{thm}

Recall that we assume that the product of paths is taken in the usual order;
that is, \index{$a*b*c*d$} $a*b*c*d\df((a*b)*c)*d$.

}

\parit{Proof.}
Suppose $a_\tau$ is a homotopy of loops at $x_0$.
Then $\bar c*a_\tau*c$ is a homotopy of loops at $x_1$.
It follows that the map 
\[\phi_c\:[a]\mapsto [\bar c*a*c]\eqlbl{eq:[f]->[hfh]}\]
is defined.
Indeed, if $a_0\sim a_1$, then $\bar c*a_0*c\sim \bar c*a_1*c$,
which means that $[\bar c*a_0*c]=[\bar c*a_1*c]$; so the right-hand side in \ref{eq:[f]->[hfh]} does not depend on the choice of loop in the homotopy class $[a]$.

Further, if $a$ and $b$ are loops based at $x_0$, then \ref{clm:homot-claims} implies that
\begin{align*}
(\bar c*a*c)*(\bar c*b*c)&\sim \bar c*a*(c*\bar c)*b*c\sim
\\
&\sim \bar c*(a*e_{x_0})*b*c\sim
\\
&\sim \bar c*(a*b)*c.
\end{align*}
In other words,
\[\phi_c([a]\cdot [b])=\phi_c[a]\cdot \phi_c[b]\quad\text{for any}\quad [a],[b]\in \pi_1(\c{X},x_0);\]
that is, $\phi_c\:\pi_1(\c{X},x_0)\to \pi_1(\c{X},x_1)$ is a homomorphism.

The same argument shows that $\phi_{\bar c}\: \pi_1(\c{X},x_1)\to \pi_1(\c{X},x_0)$ defined by
\[\phi_{\bar c}\:[d]\mapsto [c*d*\bar c]\]
is a homomorphism.

Finally note that 
\begin{align*}
c*(\bar c*a*c)*\bar c
&\sim
(c*\bar c)*a*(c*\bar c)\sim
e_{x_0}*a*e_{x_0}\sim
a
\end{align*}
for any loop $a$ based at $x_0$.
Therefore, $\phi_{\bar c}$ is a left inverse of $\phi_c$.
The same argument shows that $\phi_{\bar c}$ is also a right inverse of $\phi_c$.
It follows that $\phi_c$ is an isomorphism.
\qeds

According to the theorem, the fundamental group%
\footnote{more precisely, its \textit{isomorphism class}}
of a path-connected space does not depend on its base point.
Therefore, for a path-connected space $\c{X}$ we do not need to specify the base point of its fundamental group.
So, if we think about the group up to isomorphism, then it is safe to write $\pi_1\c{X}$ instead of $\pi_1(\c{X},x_0)$.

\begin{thm}{Exercise}\label{ex:uphi}

\begin{subthm}{ex:uphi:uphi}
 Let $f_\tau\:\c{X}\to\c{Y}$ be a homotopy.
Consider the path defined by $c(\tau)=f_\tau(x_0)$.
Show that
\[\phi_c\circ f_{0*}=f_{1*};\]
here $f_{i*}\:\pi_1(\c{X},x_0)\to\pi_1(\c{Y},f_i(x_0))$ denotes the induced homomorphism of $f_{i}$.
\end{subthm}

\begin{subthm}{ex:uphi:hom-type}
Conclude that if two path-connected topological spaces $\c{X}$ and $\c{Y}$ have the same homotopy type, then
 their fundamental groups are isomorphic.
\end{subthm}

\end{thm}


\section{Product space}

\begin{thm}{Exercise}\label{ex:hom-product}
Let $\c{X}$ and $\c{Y}$ be two path-connected topological spaces.
Choose points $x_0\in \c{X}$ and $y_0\in \c{Y}$.
Consider the projections and their induced homomorphisms:
\begin{align*}
f\:\c{X}\times\c{Y}&\to \c{X},
&
f_*\:\pi_1(\c{X}\times\c{Y},(x_0,y_0))&\to \pi_1(\c{X},x_0),
\\
g\:\c{X}\times\c{Y}&\to \c{Y},
&
g_*\:\pi_1(\c{X}\times\c{Y},(x_0,y_0))&\to \pi_1(\c{Y},y_0).
\end{align*}
Define $\phi\:\pi_1(\c{X}\times\c{Y},(x_0,y_0))\to \pi_1(\c{X},x_0)\times \pi_1(\c{Y},y_0)$ by
\[\phi\:\alpha\mapsto (f_*(\alpha),g_*(\alpha))\]
for any $\alpha\in \pi_1(\c{X}\times\c{Y},(x_0,y_0))$.
Prove the following.

\begin{subthm}{ex:hom-product:hom}
$\phi$ is a homomorphism.
\end{subthm}


\begin{subthm}{ex:hom-product:mono}
$\phi$ is a monomorphism;
that is, if $\phi(\alpha)=\phi(\beta)$ for some $\alpha,\beta\z\in \pi_1(\c{X}\times\c{Y},(x_0,y_0))$, then
$\alpha=\beta$.
\end{subthm}

\begin{subthm}{ex:hom-product:epi}
$\phi$ is an epimorphism;
that is, for any $\gamma\in \pi_1(\c{X},x_0)\z\times\pi_1(\c{Y},y_0)$ there is $\alpha\in\pi_1(\c{X}\times\c{Y},(x_0,y_0))$ such that $\phi(\alpha)=\gamma$.
\end{subthm}

Conclude that $\pi_1(\c{X}\times\c{Y})$ is isomorphic to $\pi_1(\c{X})\times\pi_1(\c{Y})$.

\end{thm}

\begin{thm}{Advanced exercise}\label{ex:sim-product}
Let $\c{X}$ be a path-connected topological space.
Consider the quotient space $\c{Y}=(\c{X}\times \c{X})/\sim$ where $\sim$ is the minimal equivalence relation such that $(x,y)\sim (y,x)$.
Show that the fundamental group of $\c{Y}$ is commutative.
\end{thm}

